\documentclass[11pt,letterpaper]{article}
\usepackage[T1]{fontenc}
\usepackage[utf8]{inputenc}
\usepackage{lmodern,microtype}
\usepackage[margin=0.78in,headheight=15pt]{geometry}
\usepackage{amsmath,amssymb,booktabs,longtable,tabularx,array,enumitem,xcolor,hyperref,fancyhdr,lastpage,needspace}
\definecolor{EngineBlue}{HTML}{183B56}
\definecolor{EngineTeal}{HTML}{117A83}
\definecolor{EngineGray}{HTML}{596671}
\definecolor{EngineLight}{HTML}{EEF4F7}
\hypersetup{colorlinks=true,urlcolor=EngineTeal,linkcolor=EngineBlue}
\setlength{\parindent}{0pt}
\setlength{\parskip}{5pt}
\emergencystretch=3em
\setlist{nosep,leftmargin=*}
\pagestyle{fancy}\fancyhf{}
\fancyhead[L]{\small\sffamily Game Engine Architecture | Engineering Suite}
\fancyhead[R]{\small\sffamily \DocType}
\fancyfoot[L]{\scriptsize Source-grounded engineering curriculum | v1.0}
\fancyfoot[R]{\scriptsize \thepage/\pageref{LastPage}}
\newcommand{\DocType}{Reference}
\newcommand{\SuiteHeading}[3]{%
\renewcommand{\DocType}{#2}%
{\small\bfseries\sffamily\color{EngineTeal}#2\hfill #3}\par\vspace{6pt}
{\LARGE\bfseries\color{EngineBlue}#1}\par\vspace{4pt}
{\color{EngineTeal}\rule{\linewidth}{1pt}}\par\vspace{7pt}}
\newcommand{\LabEvidence}[1]{\par\textbf{Evidence required:} #1\par}
\newcommand{\WriteLines}[1]{\par\begingroup\setlength{\parskip}{9pt}\count255=0\loop\ifnum\count255<#1\noindent\rule{\linewidth}{0.35pt}\par\advance\count255 by1\repeat\endgroup\smallskip}
\newcommand{\WritingBox}[1]{\noindent\fbox{\parbox[t][#1][t]{0.97\linewidth}{\mbox{}}}\par\smallskip}
\newcommand{\SourceNote}{\footnotesize\textit{Primary reading: Jason Gregory, \emph{Game Engine Architecture}, 3rd ed. (CRC Press, 2018). The narrative, exercises, laboratory designs and scoring instruments in this suite are independently authored teaching materials, not a reproduction of the textbook.}\normalsize\par}
\newcommand{\Criterion}[2]{\textbf{#1.} #2\par}
\newcommand{\FormField}[2]{\Needspace{3\baselineskip}\textbf{#1}\par\textcolor{EngineGray}{\footnotesize #2}\par\WriteLines{2}}
\newcommand{\BlankField}[1]{\Needspace{3\baselineskip}\textbf{#1}\par\WriteLines{2}}
\newcolumntype{Y}{>{\raggedright\arraybackslash}X}

\begin{document}
\SuiteHeading{Prerequisite-Based Game Engine Engineering Curriculum}{Curriculum}{24 weeks / 7 stages}
\SourceNote
\section{Entry profile and diagnostic}
Prerequisites: practical C++ (including RAII and move semantics), Python 3, unit testing, build systems, vectors/matrices, basic OS scheduling, and elementary mechanics. Read the book alongside the program; the laboratories intentionally use minimal implementations to expose contracts rather than compete with commercial engines.

\section{Dependency graph}
\begin{verbatim}
Build/debug (2) --> C++ & memory (3) --> support (6) --> resources (7)
Math (5) ------+--> rendering (11), animation (12), physics (13)
Concurrency (4) --> game loop (8) --> update scheduler (16)
Tooling (10) --> profiling & benchmarking (all implementations)
Input (9) --> game loop (8) --> gameplay object model (15,16)
Subsystems (6-14) --> architecture integration (16) --> capstone (17)
\end{verbatim}\section{24-week staged learning schedule}
\small\begin{longtable}{|p{0.235\linewidth}|p{0.235\linewidth}|p{0.235\linewidth}|p{0.235\linewidth}|}
\hline \textbf{Stage} & \textbf{Weeks} & \textbf{Reading and prerequisites} & \textbf{Mandatory assessment}\\ \hline\endhead
Foundation / tooling & 1–4 & 1–3; basic C++ and builds & Engine map; sanitizer and layout comparison\\ \hline
Math and parallelism & 5–8 & 5 then 4; threads and vectors & Transform tests; worker correctness suite\\ \hline
Low-level runtime & 9–12 & 6–8; RAII and numerical simulation & Allocator/resource tests; deterministic fixed step\\ \hline
Human interface / observability & 13–14 & 9–10; event queues and clocks & Input trace; frame profiler CSV\\ \hline
Presentation systems & 15–19 & 11–14; geometry, physics, sampling & Raster image; pose blend; collision and WAV checks\\ \hline
Gameplay foundation & 20–22 & 15–16; assets, events and update scheduling & World schema; stale-handle and event tests\\ \hline
Integration / defense & 23–24 & 17 and all evidence & Architecture package, benchmarks, review scorecard\\ \hline
\end{longtable}\normalsize
\section{Weekly operating cadence}
\begin{itemize}
\item Session A: read source sections and build concept map; list conflicting assumptions.
\item Session B: implement a baseline and one deliberately failing test.
\item Session C: instrument and rerun with a controlled parameter sweep.
\item Session D: review interface design, safety/portability and assessment evidence.
\end{itemize}
\section{Gate definitions}
\small\begin{longtable}{|p{0.47\linewidth}|p{0.47\linewidth}|}
\hline \textbf{Gate} & \textbf{Completion rule}\\ \hline\endhead
G1 Correctness & Automated tests pass, failure paths covered and reproducibility data retained.\\ \hline
G2 Timing & Median, p95, p99 and worst-case are reported with unit, frame count and environment.\\ \hline
G3 Memory & Peak footprint, ownership, reset behavior and sanitizer diagnostics documented.\\ \hline
G4 Design & Architecture decision, relevant diagrams and trade-offs recorded.\\ \hline
G5 Release & No untracked assets/licenses; no assertions about production readiness without evidence.\\ \hline
\end{longtable}\normalsize
\section{Capstone: deterministic miniature engine}
Build a loop with input event queue, generational entity storage, fixed physics, interpolation snapshot, simple raster output and synthesized audio. Include a concurrency experiment as an optional performance path; keep a single-thread reference implementation. Deliver an architecture description, dependency graph, traceability matrix, benchmark report and scored release review.


\end{document}
